% ============================================================
%  《程序员修炼之道》中英对照排版 - 样式与宏包
%  引擎：XeLaTeX（需要 fontspec / xeCJK / paracol）
% ============================================================

\usepackage{fontspec}
\usepackage{xeCJK}

% ---- 字体（macOS 系统字体）----
\setmainfont{Times New Roman}          % 英文正文
\setsansfont{Helvetica Neue}
\setmonofont{Menlo}
\setCJKmainfont{Songti SC}             % 中文正文（宋体）
\setCJKsansfont{PingFang SC}           % 中文无衬线（苹方）
\setCJKmonofont{Songti SC}

% ---- 页面 ----
\usepackage[a4paper, top=2.4cm, bottom=2.4cm, left=2.0cm, right=2.0cm]{geometry}
\usepackage{fancyhdr}
\pagestyle{fancy}
\fancyhf{}
\fancyhead[L]{\small\itshape The Pragmatic Programmer}
\fancyhead[R]{\small 程序员修炼之道}
\fancyfoot[C]{\thepage}
\renewcommand{\headrulewidth}{0.4pt}

% ---- 双语并排 ----
\usepackage{paracol}
\columnratio{0.5}
\setlength{\columnsep}{1.8em}
\setlength{\parindent}{0pt}
\setlength{\parskip}{0.55em}
\linespread{1.12}

% 对照宏：\src{英文} / \tgt{中文}
% 关键：每对只在 \src 起始处同步一次（* 号），\tgt 切换时不再同步，
% 否则两次同步会互相追赶造成逐段累积错位。
\newcommand{\src}[1]{\switchcolumn[0]*\noindent #1\par}
\newcommand{\tgt}[1]{\switchcolumn[1]\noindent #1\par}

% 不必成对时，直接开环境
\newenvironment{bitext}{\begin{paracol}{2}\raggedright}{\end{paracol}}

% ---- 标题层级 ----
\usepackage{titlesec}
\usepackage{xcolor}
\definecolor{ppblue}{HTML}{1F4E79}
\definecolor{ppgray}{HTML}{555555}
\definecolor{tipbg}{HTML}{F2F6FA}

\titleformat{\chapter}[display]
  {\normalfont\huge\bfseries\color{ppblue}}
  {\filright\normalsize\sffamily CHAPTER \thechapter}
  {1ex}
  {\Huge}
\titlespacing*{\chapter}{0pt}{-10pt}{24pt}

\titleformat{\section}
  {\normalfont\Large\bfseries\color{ppblue}}{\thesection}{0.8em}{}
\titleformat{\subsection}
  {\normalfont\large\bfseries\color{ppblue}}{\thesubsection}{0.7em}{}

% ---- 列表 ----
\usepackage{enumitem}
\setlist{nosep, leftmargin=1.4em}

% ---- Tip 提示框 ----
\usepackage[most]{tcolorbox}
\newtcolorbox{tipbox}[1][]{
  colback=tipbg, colframe=ppblue, boxrule=0.6pt, arc=2pt,
  left=6pt, right=6pt, top=5pt, bottom=5pt, #1
}
\newcounter{tipcounter}
% \tip{<全书编号>}{<英文标题>}{<说明>}
\newcommand{\tip}[3]{%
  \begin{tipbox}
    {\footnotesize\sffamily\bfseries\color{ppblue}TIP #1}\par
    \vspace{2pt}{\bfseries #2}\par
    \vspace{2pt}{\itshape\small\color{ppgray} #3}
  \end{tipbox}%
}

% ---- 杂项 ----
\usepackage{hyperref}
\hypersetup{colorlinks=true, linkcolor=ppblue, urlcolor=ppblue}
\usepackage{booktabs}
\usepackage{longtable}

\newcommand{\chapterplaceholder}[2]{%
  \chapter{#1}
  \begin{center}
  \begin{minipage}{0.86\textwidth}
  \itshape\small\color{ppgray}
  #2
  \end{minipage}
  \end{center}
}

% ---- 双语导读版式 ----
% 小节标题（英文 + 中文）
\newcommand{\sechead}[2]{%
  \section*{#1\\[0.2em]{\large\normalfont #2}}%
  \addcontentsline{toc}{section}{#1}%
}

% 核心论点框
\newtcolorbox{ideabox}[1][核心论点]{%
  enhanced, colback=tipbg, colframe=ppblue!70, boxrule=0.6pt, arc=2pt,
  left=8pt, right=8pt, top=6pt, bottom=6pt,
  colbacktitle=ppblue, coltitle=white,
  fonttitle=\bfseries\sffamily\small,
  title={#1},
}

% 解读（中文评论）
\newcommand{\commentary}[1]{%
  \par\medskip
  \noindent{\small\bfseries\color{ppblue}解读\ }\ignorespaces
  {\small #1}\par
}

% 实践清单
\newenvironment{actions}{%
  \par\medskip\noindent{\small\bfseries\color{ppblue}实践清单}\par
  \begin{itemize}[leftmargin=1.4em,itemsep=1pt,topsep=2pt]
    \small
}{%
  \end{itemize}\par
}
